\begin{table}[!htbp]
\centering
\caption{\textbf{Player 1: Counterpart Proximity within Stated Reason Category}}
\label{tab:player1_sp_by_category}
\begin{tabular}{lccc}
\toprule
& Moral & Self-interest & Mutual Benefit/Cooperation \\
\midrule
\multicolumn{4}{l}{\emph{Panel A: Distribution of counterpart proximity (column shares)}} \\ % chktex 13
Low social proximity & 49.4\% & 91.9\% & 59.4\% \\
High social proximity & 50.6\% & 8.1\% & 40.6\% \\
\midrule
\multicolumn{4}{l}{\emph{Panel B: Mean action, share of the budget}} \\ % chktex 13
Low social proximity & 44.3\% & 23.6\% & 57.0\% \\
High social proximity & 47.0\% & 34.6\% & 63.2\% \\
\bottomrule
\end{tabular}
\begin{flushleft}
\footnotesize Notes: Pooled across all games and conditions for Player 1, classified sample. Panel A reports the distribution of counterpart proximity within each stated reason category (columns sum to 100\%). Panel B reports the mean action (share of the budget transferred, offered, or sent) in each proximity-by-category cell. High social proximity: the stated reason refers to the counterpart as at least an anonymous peer (anonymous peer, teammate or coworker, friend); low: no mention of the counterpart or an abstract stranger.
\end{flushleft}
\end{table}
